\documentclass[12pt, a4paper]{article}
\usepackage[utf8]{inputenc}
\usepackage[T1]{fontenc}
\usepackage[margin=2.5cm]{geometry}
\usepackage{amsmath}
\usepackage{amssymb}
\usepackage{enumitem}
\usepackage{xcolor}
\usepackage{parskip}

\definecolor{truegreen}{RGB}{39, 80, 10}
\definecolor{falsered}{RGB}{121, 31, 31}
\definecolor{darkblue}{RGB}{10, 40, 120}

\newcommand{\T}{\textcolor{truegreen}{\textbf{[TRUE]}}}
\newcommand{\F}{\textcolor{falsered}{\textbf{[FALSE]}}}
\newcommand{\expl}[1]{\textcolor{darkblue}{\small #1}}

\begin{document}

\begin{center}
  {\LARGE \textbf{Multivariate Normal Distribution (MVN)}}\\[0.3em]
  {\large Exam Questions --- Answers \& Explanations}
\end{center}
\vspace{1em}

\textbf{Key rules to memorise:}
\begin{itemize}[itemsep=4pt]
  \item \textbf{Linear transformations of a MVN are always MVN.} There are no exceptions.
  \item \textbf{Marginals of a MVN are univariate normal.} But the converse is false: normal marginals do \textit{not} imply a joint MVN.
  \item \textbf{MVN is characterised by every linear combination being normal.} If \textit{all} linear combinations $\mathbf{a}^\top \mathbf{X}$ are univariate normal, then $\mathbf{X}$ is MVN.
  \item \textbf{For jointly Gaussian variables: uncorrelated $\Leftrightarrow$ independent.} This equivalence does \textit{not} hold in general.
  \item \textbf{In general: uncorrelated $\not\Rightarrow$ independent}, and \textbf{dependent $\not\Rightarrow$ correlated} (nonlinear dependence can exist with zero correlation).
  \item The covariance matrix $\boldsymbol{\Sigma}$ is \textbf{positive semi-definite}, so $\det(\boldsymbol{\Sigma}) \geq 0$.
  \item A MVN is \textbf{fully characterised by $\boldsymbol{\mu}$ and $\boldsymbol{\Sigma}$} alone.
  \item \textbf{Squared Mahalanobis distance} $(\mathbf{X}_i - \boldsymbol{\mu})^\top \boldsymbol{\Sigma}^{-1} (\mathbf{X}_i - \boldsymbol{\mu}) \sim \chi^2(p)$.
  \item \textbf{Wishart distribution} is the multivariate generalisation of the chi-square. It is additive (reproductive) and reduces to a scaled $\chi^2$ when $p=1$.
\end{itemize}

\vspace{1.5em}

% -----------------------------------------------------------------------
\section*{Q1. Suppose $\mathbf{X}$ follows a MVN with covariance matrix $\boldsymbol{\Sigma}$. Which statements are correct?}

\begin{enumerate}[label=(\arabic*)]
  \item \F\ In some cases, it is possible to find a linear transformation of $\mathbf{X}$ that does not result in another MVN.
  \item \T\ If $\mathbf{X}$ is normally distributed, all of its marginal distributions are also normal.
  \item \T\ If the components of $\mathbf{X}$ are uncorrelated, then they are independent.
  \item \T\ The determinant of $\boldsymbol{\Sigma}$ is greater than or equal to zero.
\end{enumerate}

\expl{
\textbf{(1)} False. Any linear transformation $\mathbf{A}\mathbf{X} + \mathbf{b}$ of a MVN is always MVN. This is one of the defining closure properties; there are no exceptions.\\
\textbf{(2)} True. Each individual component $X_i \sim \mathcal{N}(\mu_i, \sigma_i^2)$. Marginals of a MVN are always univariate normal.\\
\textbf{(3)} True --- but \textit{only because} $\mathbf{X}$ is jointly Gaussian. Zero covariance (uncorrelatedness) implies independence when the joint distribution is Gaussian. This does not hold in general.\\
\textbf{(4)} True. $\boldsymbol{\Sigma}$ is positive semi-definite by definition, so all eigenvalues are $\geq 0$, which implies $\det(\boldsymbol{\Sigma}) \geq 0$.
}

\vspace{1.5em}

% -----------------------------------------------------------------------
\section*{Q2. Let $\mathbf{X}$ be a random vector with values in $\mathbb{R}^p$. Indicate the true statements:}

\begin{enumerate}[label=(\arabic*)]
  \item \T\ If $\mathbf{X}$ is multivariate Gaussian, the marginal distributions of $\mathbf{X}$ are univariate Gaussian.
  \item \T\ If $\mathbf{X}$ is multivariate Gaussian, any linear combination of the coordinates of $\mathbf{X}$ is univariate Gaussian.
  \item \F\ $\mathbf{X}$ is multivariate Gaussian if and only if the marginal distributions of $\mathbf{X}$ are univariate Gaussian.
  \item \F\ If one can find a linear combination of the coordinates of $\mathbf{X}$ that is univariate Gaussian, then $\mathbf{X}$ is multivariate Gaussian.
\end{enumerate}

\expl{
\textbf{(1)} True. Projecting a MVN onto any coordinate axis yields a univariate normal.\\
\textbf{(2)} True. This is actually a standard \textit{definition} of the MVN: $\mathbf{X}$ is MVN if and only if \textit{every} linear combination $\mathbf{a}^\top \mathbf{X}$ is univariate normal.\\
\textbf{(3)} False. Normal marginals do \textit{not} guarantee a MVN joint distribution. A classic counterexample: one can construct $(X, Y)$ where $X$ and $Y$ are each marginally normal, but $(X, Y)$ is not jointly normal.\\
\textbf{(4)} False. Finding \textit{one} linear combination that happens to be Gaussian is not sufficient. \textit{All} linear combinations must be Gaussian for $\mathbf{X}$ to be MVN.
}

\vspace{1.5em}

% -----------------------------------------------------------------------
\section*{Q3. Indicate the true statements about uncorrelatedness and independence:}

\begin{enumerate}[label=(\arabic*)]
  \item \T\ Two random variables whose joint distribution is Gaussian are independent if and only if they are uncorrelated.
  \item \T\ If two random variables whose joint distribution is Gaussian are uncorrelated, then they are independent.
  \item \F\ If two random variables are uncorrelated, then they are independent.
  \item \F\ If two random variables are dependent, then they are correlated.
\end{enumerate}

\expl{
\textbf{(1)} True. For jointly Gaussian variables, uncorrelated $\Leftrightarrow$ independent. Both directions hold.\\
\textbf{(2)} True. This is the key special property of the Gaussian: $\mathrm{Cov}(X, Y) = 0$ implies independence, given that the joint distribution is Gaussian.\\
\textbf{(3)} False in general. Uncorrelated means only zero \textit{linear} correlation. Nonlinear dependence can still exist. The Gaussian case is a special exception.\\
\textbf{(4)} False. Dependence does not imply correlation. Classic example: let $X \sim \mathcal{N}(0,1)$ and $Y = X^2$. Then $Y$ is completely determined by $X$ (dependent), yet $\mathrm{Cov}(X, Y) = 0$ (uncorrelated).
}

\vspace{1.5em}

% -----------------------------------------------------------------------
\section*{Q4. Let $\mathbf{X}$ be a $p$-dimensional random vector following a MVN. Which statements are true?}

\begin{enumerate}[label=(\arabic*)]
  \item \T\ If all pairwise correlations between the components of $\mathbf{X}$ are zero, they are independent.
  \item \T\ Any linear combination of the components of $\mathbf{X}$ is normally distributed.
  \item \T\ The marginal distributions of $\mathbf{X}$ are univariate normal.
  \item \T\ The joint distribution of $\mathbf{X}$ is completely characterised by its mean vector and covariance matrix.
\end{enumerate}

\expl{
\textbf{(1)} True. Zero pairwise correlations means all off-diagonal entries of $\boldsymbol{\Sigma}$ are zero, so $\boldsymbol{\Sigma}$ is diagonal. For a MVN, a diagonal covariance matrix implies full independence.\\
\textbf{(2)} True. For any $\mathbf{a} \in \mathbb{R}^p$: $\mathbf{a}^\top \mathbf{X} \sim \mathcal{N}(\mathbf{a}^\top \boldsymbol{\mu},\; \mathbf{a}^\top \boldsymbol{\Sigma}\, \mathbf{a})$.\\
\textbf{(3)} True. Each $X_i \sim \mathcal{N}(\mu_i, \sigma_i^2)$. Marginals of a MVN are always univariate normal.\\
\textbf{(4)} True. The MVN is fully determined by its mean vector $\boldsymbol{\mu}$ and covariance matrix $\boldsymbol{\Sigma}$. No additional parameters are needed.
}

\vspace{1.5em}

% -----------------------------------------------------------------------
\section*{Q5. Let $\mathbf{X}_1, \ldots, \mathbf{X}_n$ be iid from a $p$-dimensional MVN. Which statements are correct?}

\begin{enumerate}[label=(\arabic*)]
  \item \T\ The squared Mahalanobis distances of the observations from the true mean follow a chi-square distribution.
  \item \F\ If $n > p$, then the sample covariance matrix is certainly singular.
  \item \F\ The total variance of the sample is the product of the squared Euclidean distances from each observation to the sample mean.
  \item \T\ The squared Mahalanobis distance of the sample mean from the true mean follows a chi-square distribution.
\end{enumerate}

\expl{
\textbf{(1)} True. $(\mathbf{X}_i - \boldsymbol{\mu})^\top \boldsymbol{\Sigma}^{-1} (\mathbf{X}_i - \boldsymbol{\mu}) \sim \chi^2(p)$. This follows directly from standardising a MVN.\\
\textbf{(2)} False --- it is the \textit{opposite}. When $n > p$, the sample covariance matrix is almost surely \textit{non-singular} (full rank). Singularity is guaranteed when $n \leq p$, because there are too few observations to span $\mathbb{R}^p$.\\
\textbf{(3)} False. Total variance $= \mathrm{tr}(\mathbf{S}) = \sum_{j=1}^p s_j^2$, which is a \textit{sum}. The product of eigenvalues gives the determinant (generalised variance), which is a completely different quantity.\\
\textbf{(4)} True. Since $\bar{\mathbf{X}} \sim \mathcal{N}(\boldsymbol{\mu},\, \boldsymbol{\Sigma}/n)$, it follows that $n(\bar{\mathbf{X}} - \boldsymbol{\mu})^\top \boldsymbol{\Sigma}^{-1} (\bar{\mathbf{X}} - \boldsymbol{\mu}) \sim \chi^2(p)$.
}

\vspace{1.5em}

% -----------------------------------------------------------------------
\section*{Q6. Let $\mathbf{X}_1, \ldots, \mathbf{X}_n$ be iid from a $p$-dimensional MVN with $n > p$. Which statements about the Wishart distribution are true?}

\begin{enumerate}[label=(\arabic*)]
  \item \F\ The sample covariance matrix follows a Wishart distribution.
  \item \T\ If $p = 1$, the Wishart distribution is equivalent to a scaled chi-square distribution.
  \item \F\ The Wishart distribution is always defined, even when the covariance matrix is singular.
  \item \T\ The sum of two independent random matrices, each with a Wishart distribution with the same covariance matrix as parameter, has a Wishart distribution.
\end{enumerate}

\expl{
\textbf{(1)} False --- not quite. It is the \textit{scatter matrix} $(n-1)\mathbf{S} = \sum_{i=1}^n (\mathbf{X}_i - \bar{\mathbf{X}})(\mathbf{X}_i - \bar{\mathbf{X}})^\top$ that follows $\mathcal{W}_p(n-1,\, \boldsymbol{\Sigma})$. The sample covariance $\mathbf{S}$ itself is a scaled Wishart, not a Wishart directly.\\
\textbf{(2)} True. When $p = 1$: $\mathcal{W}_1(k, \sigma^2) \equiv \sigma^2 \chi^2(k)$. The Wishart generalises the chi-square distribution to multiple dimensions.\\
\textbf{(3)} False. The standard Wishart distribution requires $\boldsymbol{\Sigma}$ to be positive \textit{definite} (non-singular). Generalisations exist for singular cases, but the standard definition excludes them.\\
\textbf{(4)} True. This is the \textit{reproductive (additive) property} of the Wishart: if $\mathbf{A} \sim \mathcal{W}_p(m, \boldsymbol{\Sigma})$ and $\mathbf{B} \sim \mathcal{W}_p(n, \boldsymbol{\Sigma})$ independently, then $\mathbf{A} + \mathbf{B} \sim \mathcal{W}_p(m+n, \boldsymbol{\Sigma})$.
}


%%% -----------------------------------------------------------------------
%%% Bu bloğu mvn_exam_notes.tex dosyasının \end{document} satırından
%%% ÖNCE yapıştırın.
%%% -----------------------------------------------------------------------

\vspace{2em}
\begin{center}
  {\large \textbf{Additional Questions --- Lecture 7 \& 8 Material}}
\end{center}
\vspace{0.5em}

% -----------------------------------------------------------------------
\section*{Q7. Suppose $\mathbf{x} \sim N_p(\boldsymbol{\mu}, \boldsymbol{\Sigma})$. Which statements about the density and geometry are correct?}

\begin{enumerate}[label=(\arabic*)]
  \item \T\ The level sets of the density $f_{\mathbf{x}}$ are ellipsoids centred at $\boldsymbol{\mu}$.
  \item \T\ The eigenvectors of $\boldsymbol{\Sigma}$ give the axes of these ellipsoids.
  \item \F\ If $\boldsymbol{\Sigma} = \mathbf{I}$, the level sets are ellipsoids with unequal axes.
  \item \T\ Larger eigenvalues $\lambda_i$ correspond to directions along which the ellipsoid is more stretched.
\end{enumerate}

\expl{
\textbf{(1)} True. $f_{\mathbf{x}}(\mathbf{t}) = c \Leftrightarrow (\mathbf{t} - \boldsymbol{\mu})^\top \boldsymbol{\Sigma}^{-1}(\mathbf{t} - \boldsymbol{\mu}) = \text{const}$, which defines an ellipsoid centred at $\boldsymbol{\mu}$.\\
\textbf{(2)} True. The spectral decomposition $\boldsymbol{\Sigma} = \sum_i \lambda_i \mathbf{e}_i \mathbf{e}_i^\top$ shows the eigenvectors $\mathbf{e}_i$ are the principal axes and $\sqrt{\lambda_i}$ are the half-lengths.\\
\textbf{(3)} False. When $\boldsymbol{\Sigma} = \mathbf{I}$, all eigenvalues equal 1 and all directions are equally scaled, so the level sets are perfect \textit{spheres}, not stretched ellipsoids.\\
\textbf{(4)} True. A large $\lambda_i$ means high variance in direction $\mathbf{e}_i$, which corresponds to greater spread — and therefore greater half-length $\sqrt{\lambda_i}$ — along that axis.
}

\vspace{1.5em}

% -----------------------------------------------------------------------
\section*{Q8. Let $\mathbf{x} \sim N_p(\boldsymbol{\mu}, \boldsymbol{\Sigma})$ with block partition $\mathbf{x} = (\mathbf{x}_1^\top, \mathbf{x}_2^\top)^\top$. Which statements are correct?}

\begin{enumerate}[label=(\arabic*)]
  \item \T\ The marginal distribution of $\mathbf{x}_1$ is $N_q(\boldsymbol{\mu}_1, \boldsymbol{\Sigma}_{11})$.
  \item \T\ If $\boldsymbol{\Sigma}_{12} = \mathbf{0}$, then $\mathbf{x}_1$ and $\mathbf{x}_2$ are independent.
  \item \F\ If $\boldsymbol{\Sigma}_{12} = \mathbf{0}$, then $\mathbf{x}_1$ and $\mathbf{x}_2$ are uncorrelated but not necessarily independent.
  \item \T\ The conditional covariance $\mathrm{Cov}(\mathbf{x}_1 \mid \mathbf{X}_2 = \mathbf{x}_2) = \boldsymbol{\Sigma}_{11} - \boldsymbol{\Sigma}_{12}\boldsymbol{\Sigma}_{22}^{-1}\boldsymbol{\Sigma}_{21}$ does not depend on the observed value $\mathbf{x}_2$.
\end{enumerate}

\expl{
\textbf{(1)} True. Apply $A = [I_q \mid \mathbf{0}]$; then $A\mathbf{x} = \mathbf{x}_1$ and $A\boldsymbol{\Sigma} A^\top = \boldsymbol{\Sigma}_{11}$.\\
\textbf{(2)} True. This is the special Gaussian property: for jointly Gaussian blocks, zero cross-covariance ($\boldsymbol{\Sigma}_{12} = \mathbf{0}$) implies full statistical independence (the joint density factors).\\
\textbf{(3)} False. For jointly Gaussian variables the two statements are equivalent: $\boldsymbol{\Sigma}_{12} = \mathbf{0} \Leftrightarrow \mathbf{x}_1 \perp\!\!\!\perp \mathbf{x}_2$. The distinction between uncorrelated and independent only matters outside the Gaussian setting.\\
\textbf{(4)} True. The Schur complement $\boldsymbol{\Sigma}_{11} - \boldsymbol{\Sigma}_{12}\boldsymbol{\Sigma}_{22}^{-1}\boldsymbol{\Sigma}_{21}$ is a fixed matrix — it depends on the population parameters, not on the specific realisation $\mathbf{x}_2$. Only the conditional \textit{mean} shifts with $\mathbf{x}_2$.
}

\vspace{1.5em}

% -----------------------------------------------------------------------
\section*{Q9. Let $(Y, X)^\top \sim N_2$ with correlation $\rho = \mathrm{corr}(X, Y)$. Which statements are true?}

\begin{enumerate}[label=(\arabic*)]
  \item \T\ The conditional mean $\mathbb{E}[Y \mid X = x]$ is a linear function of $x$.
  \item \T\ In standardized coordinates, the regression line has slope exactly $\rho$.
  \item \F\ If $\rho = 0.8$, then an extremely high $X$ value predicts an equally extreme $Y$ value on average.
  \item \T\ The regression fallacy consists of giving a causal explanation to a purely mathematical phenomenon.
\end{enumerate}

\expl{
\textbf{(1)} True. From the conditional theorem: $\mathbb{E}[Y \mid X = x] = \mu_Y + \frac{\sigma_{XY}}{\sigma_{XX}}(x - \mu_X)$, which is linear in $x$.\\
\textbf{(2)} True. In standardized coordinates $y_{\mathrm{std}} = \rho \cdot x_{\mathrm{std}}$. The slope is exactly $\rho$.\\
\textbf{(3)} False. Because $|\rho| \leq 1$, the regression line is always \textit{flatter} than the 45° line. Even with $\rho = 0.8$, an extreme $X$ predicts a $Y$ that is closer to its mean (in standard deviation units) than $X$ is to its own mean. This is regression toward the mean.\\
\textbf{(4)} True. Regression toward the mean is a mathematical inevitability from $|\rho| \leq 1$. Interpreting it as evidence of a biological, economic, or causal mechanism is the fallacy. Galton's tall-fathers observation and the cited economics Nobel mistakes are canonical examples.
}

\vspace{1.5em}

% -----------------------------------------------------------------------
\section*{Q10. Let $X_1, \ldots, X_n \overset{\mathrm{iid}}{\sim} N_p(\boldsymbol{\mu}, \boldsymbol{\Sigma})$. Which statements about the MLE are correct?}

\begin{enumerate}[label=(\arabic*)]
  \item \T\ The MLE of $\boldsymbol{\mu}$ is the sample mean $\bar{\mathbf{X}}$, which is unbiased.
  \item \F\ The MLE of $\boldsymbol{\Sigma}$ is the sample covariance matrix $\mathbf{S}$ (dividing by $n-1$), which is unbiased.
  \item \T\ The MLE of $\boldsymbol{\Sigma}$ divides by $n$, not $n-1$, and is therefore biased.
  \item \T\ By the invariance property of the MLE, the MLE of $\det(\boldsymbol{\Sigma})$ is $\det(\hat{\boldsymbol{\Sigma}})$.
\end{enumerate}

\expl{
\textbf{(1)} True. $\hat{\boldsymbol{\mu}} = \bar{\mathbf{X}} = \frac{1}{n}\sum X_i$ and $\mathbb{E}[\bar{\mathbf{X}}] = \boldsymbol{\mu}$.\\
\textbf{(2)} False. The MLE of $\boldsymbol{\Sigma}$ divides by $n$: $\hat{\boldsymbol{\Sigma}} = \frac{1}{n}\sum (X_i - \bar{\mathbf{X}})(X_i - \bar{\mathbf{X}})^\top = \frac{n-1}{n}\mathbf{S}$. This is biased: $\mathbb{E}[\hat{\boldsymbol{\Sigma}}] = \frac{n-1}{n}\boldsymbol{\Sigma} \neq \boldsymbol{\Sigma}$. The unbiased estimator $\mathbf{S}$ is \textit{not} the MLE.\\
\textbf{(3)} True. Bessel's correction ($n-1$ denominator) corrects for the fact that deviations $(X_i - \bar{\mathbf{X}})$ are constrained to sum to zero, making them systematically smaller. The MLE ignores this correction and is therefore biased downward.\\
\textbf{(4)} True. The invariance property states: if $\hat{\theta}$ is the MLE of $\theta$, then $h(\hat{\theta})$ is the MLE of $h(\theta)$ for any function $h$. No re-derivation is needed — just apply $h = \det(\cdot)$ to $\hat{\boldsymbol{\Sigma}}$.
}

\vspace{1.5em}

% -----------------------------------------------------------------------
\section*{Q11. Let $X_1, \ldots, X_n \overset{\mathrm{iid}}{\sim} N_p(\boldsymbol{\mu}, \boldsymbol{\Sigma})$ with $n > p$. Which statements about the distribution of $\bar{\mathbf{X}}$ and $\mathbf{S}$ are true?}

\begin{enumerate}[label=(\arabic*)]
  \item \T\ $\bar{\mathbf{X}} \sim N_p\!\left(\boldsymbol{\mu},\, \frac{1}{n}\boldsymbol{\Sigma}\right)$ exactly for any $n \geq 1$.
  \item \T\ $\bar{\mathbf{X}}$ and $\mathbf{S}$ are independent.
  \item \F\ $(n-1)\mathbf{S} \sim W_p(n, \boldsymbol{\Sigma})$, i.e., with $n$ degrees of freedom.
  \item \T\ $n(\bar{\mathbf{X}} - \boldsymbol{\mu})^\top \boldsymbol{\Sigma}^{-1}(\bar{\mathbf{X}} - \boldsymbol{\mu}) \sim \chi^2(p)$.
\end{enumerate}

\expl{
\textbf{(1)} True. The Gaussian is closed under sums and scalar multiplication. The result is exact (not an approximation) for all $n \geq 1$.\\
\textbf{(2)} True. This is one of the most important structural facts in multivariate statistics. For Gaussian data, the location of the centre ($\bar{\mathbf{X}}$) carries no information about the spread ($\mathbf{S}$) and vice versa. This independence is what makes Hotelling's $T^2$ test possible.\\
\textbf{(3)} False. The degrees of freedom are $n-1$, not $n$. The constraint $\sum(X_i - \bar{\mathbf{X}}) = \mathbf{0}$ removes one degree of freedom. The correct result is $(n-1)\mathbf{S} \sim W_p(n-1, \boldsymbol{\Sigma})$.\\
\textbf{(4)} True. Since $\bar{\mathbf{X}} \sim N_p(\boldsymbol{\mu}, \boldsymbol{\Sigma}/n)$, standardizing gives $\sqrt{n}\,\boldsymbol{\Sigma}^{-1/2}(\bar{\mathbf{X}} - \boldsymbol{\mu}) \sim N_p(\mathbf{0}, \mathbf{I})$, and the squared norm of a standard $p$-variate normal is $\chi^2(p)$.
}

\vspace{1.5em}

% -----------------------------------------------------------------------
\section*{Q12. Which of the following statements about the Law of Large Numbers (LLN) and the Central Limit Theorem (CLT) are correct?}

\begin{enumerate}[label=(\arabic*)]
  \item \T\ The LLN holds for any distribution with finite mean and covariance --- Gaussianity is not required.
  \item \F\ The CLT states that $\bar{\mathbf{X}}_n - \boldsymbol{\mu}$ converges in distribution to $N_p(\mathbf{0}, \boldsymbol{\Sigma})$.
  \item \T\ The typical magnitude of the estimation error satisfies $\|\bar{\mathbf{X}}_n - \boldsymbol{\mu}\| \sim C/\sqrt{n}$, so halving the error requires four times as much data.
  \item \F\ For large $n$, the sample covariance $\mathbf{S}$ follows an approximate Wishart distribution regardless of the underlying distribution of the data.
\end{enumerate}

\expl{
\textbf{(1)} True. Both the weak and strong LLN require only that the data is iid with finite mean (and finite covariance for the multivariate case). The distribution can be Bernoulli, Poisson, uniform — anything.\\
\textbf{(2)} False. $\bar{\mathbf{X}}_n - \boldsymbol{\mu} \to \mathbf{0}$, so its distribution collapses to a point mass. The correct statement is that the \textit{rescaled} error $\sqrt{n}(\bar{\mathbf{X}}_n - \boldsymbol{\mu}) \xrightarrow{d} N_p(\mathbf{0}, \boldsymbol{\Sigma})$. The $\sqrt{n}$ factor is essential --- it zooms in at exactly the right rate.\\
\textbf{(3)} True. Since $\mathrm{Var}(\bar{\mathbf{X}}_n) = \boldsymbol{\Sigma}/n$, the standard deviation shrinks at rate $1/\sqrt{n}$. This is the fundamental law of statistical estimation: cutting the error by a factor of $k$ requires $k^2$ times the data.\\
\textbf{(4)} False. The Wishart distribution of $(n-1)\mathbf{S}$ is an \textit{exact} result that requires Gaussian data. For non-Gaussian data, no simple closed-form distribution for $\mathbf{S}$ is available, even asymptotically. What the CLT guarantees is only the approximate normality of $\bar{\mathbf{X}}_n$, not any distributional result for $\mathbf{S}$.
}

\vspace{2em}
\hrule
\vspace{1.5em}

% -----------------------------------------------------------------------
\section*{Exam Reference Sheet --- Key Facts to Memorise}

\subsection*{MVN: Core Properties}
\begin{itemize}[itemsep=5pt]
  \item \textbf{Closure under linear maps.} $\mathbf{x} \sim N_p(\boldsymbol{\mu}, \boldsymbol{\Sigma})$ and $A$ a $q \times p$ matrix $\Rightarrow A\mathbf{x} \sim N_q(A\boldsymbol{\mu},\, A\boldsymbol{\Sigma} A^\top)$. No exceptions.
  \item \textbf{Characterisation theorem.} $\mathbf{x} \sim N_p \Leftrightarrow \mathbf{a}^\top \mathbf{x} \sim N_1$ for \textit{every} $\mathbf{a} \in \mathbb{R}^p$. This is both the definition and the main proof tool.
  \item \textbf{Normal marginals $\not\Rightarrow$ joint normal.} The $\Leftarrow$ direction fails. You can construct distributions with all Gaussian marginals that are not jointly Gaussian.
  \item \textbf{Gaussian special property.} For jointly Gaussian variables: uncorrelated $\Leftrightarrow$ independent ($\boldsymbol{\Sigma}_{12} = \mathbf{0} \Leftrightarrow \mathbf{x}_1 \perp\!\!\!\perp \mathbf{x}_2$). This equivalence does \textbf{not} hold in general.
  \item \textbf{Covariance matrix.} $\boldsymbol{\Sigma}$ is positive semi-definite $\Rightarrow$ all eigenvalues $\geq 0 \Rightarrow \det(\boldsymbol{\Sigma}) \geq 0$.
  \item \textbf{Full characterisation.} A MVN is completely determined by $\boldsymbol{\mu}$ and $\boldsymbol{\Sigma}$ alone.
\end{itemize}

\subsection*{Mahalanobis Distance and Confidence Ellipsoids}
\begin{itemize}[itemsep=5pt]
  \item \textbf{Squared Mahalanobis distance.} $W = (\mathbf{x} - \boldsymbol{\mu})^\top \boldsymbol{\Sigma}^{-1} (\mathbf{x} - \boldsymbol{\mu}) \sim \chi^2(p)$.
  \item \textbf{Proof strategy.} Let $\mathbf{z} = \boldsymbol{\Sigma}^{-1/2}(\mathbf{x}-\boldsymbol{\mu}) \sim N_p(\mathbf{0}, \mathbf{I})$; then $W = \mathbf{z}^\top\mathbf{z} = \sum z_i^2 \sim \chi^2(p)$.
  \item \textbf{Confidence ellipsoid.} $P\!\left[(\mathbf{x}-\boldsymbol{\mu})^\top\boldsymbol{\Sigma}^{-1}(\mathbf{x}-\boldsymbol{\mu}) \leq \chi^2_{1-\alpha}(p)\right] = 1 - \alpha$. The inverse problem: after observing $\mathbf{x}$, the ellipsoid centred at $\mathbf{x}$ covers the unknown $\boldsymbol{\mu}$ with probability $1-\alpha$.
  \item \textbf{Sample mean version.} $n(\bar{\mathbf{X}} - \boldsymbol{\mu})^\top \boldsymbol{\Sigma}^{-1}(\bar{\mathbf{X}} - \boldsymbol{\mu}) \sim \chi^2(p)$ exactly (for Gaussian data).
\end{itemize}

\subsection*{Conditional Distribution}
\begin{itemize}[itemsep=5pt]
  \item $\mathbf{x}_1 \mid \mathbf{X}_2 = \mathbf{x}_2 \sim N_q\!\left(\boldsymbol{\mu}_1 + \boldsymbol{\Sigma}_{12}\boldsymbol{\Sigma}_{22}^{-1}(\mathbf{x}_2 - \boldsymbol{\mu}_2),\; \boldsymbol{\Sigma}_{11} - \boldsymbol{\Sigma}_{12}\boldsymbol{\Sigma}_{22}^{-1}\boldsymbol{\Sigma}_{21}\right)$.
  \item \textbf{Conditional mean} shifts linearly with $\mathbf{x}_2$; \textbf{conditional covariance} (Schur complement) is fixed --- it does \textit{not} depend on the observed value $\mathbf{x}_2$.
  \item Observing $\mathbf{x}_2$ \textbf{never increases} uncertainty: $\boldsymbol{\Sigma}_{11} - \boldsymbol{\Sigma}_{12}\boldsymbol{\Sigma}_{22}^{-1}\boldsymbol{\Sigma}_{21} \leq \boldsymbol{\Sigma}_{11}$.
  \item \textbf{Regression line (bivariate).} $\mathbb{E}[Y \mid X = x] = \mu_Y + \rho\frac{\sqrt{\sigma_{YY}}}{\sqrt{\sigma_{XX}}}(x - \mu_X)$. In standardised coordinates: slope $= \rho$.
  \item \textbf{Regression fallacy.} Because $|\rho| \leq 1$, extreme $X$ predicts a less extreme $Y$ (in SD units). This is mathematical, \textbf{not causal}.
\end{itemize}

\subsection*{MLE for the Gaussian}
\begin{itemize}[itemsep=5pt]
  \item $\hat{\boldsymbol{\mu}} = \bar{\mathbf{X}}$ (MLE and unbiased). \quad $\hat{\boldsymbol{\Sigma}} = \frac{1}{n}\sum(X_i - \bar{\mathbf{X}})(X_i-\bar{\mathbf{X}})^\top = \frac{n-1}{n}\mathbf{S}$ (MLE, \textbf{biased}).
  \item $\mathbf{S} = \frac{1}{n-1}\sum(X_i-\bar{\mathbf{X}})(X_i-\bar{\mathbf{X}})^\top$ is \textbf{unbiased} but \textbf{not} the MLE.
  \item \textbf{Invariance property.} The MLE of $h(\theta)$ is $h(\hat{\theta})$. Never re-derive --- apply the function directly to $\hat{\boldsymbol{\Sigma}}$.
  \item \textbf{Why biased?} $\sum(X_i - \bar{\mathbf{X}}) = \mathbf{0}$ always (one constraint) $\Rightarrow$ one degree of freedom lost $\Rightarrow$ deviations systematically underestimated.
\end{itemize}

\subsection*{Wishart Distribution}
\begin{itemize}[itemsep=5pt]
  \item \textbf{Definition.} $W = \sum_{j=1}^m Z_j Z_j^\top$, $Z_j \overset{\mathrm{iid}}{\sim} N_p(\mathbf{0}, \boldsymbol{\Sigma})$ $\Rightarrow$ $W \sim W_p(m, \boldsymbol{\Sigma})$. Slogan: \textbf{Wishart is the matrix $\chi^2$}.
  \item \textbf{Connection to $\mathbf{S}$.} $(n-1)\mathbf{S} \sim W_p(n-1, \boldsymbol{\Sigma})$, \textbf{not} $n$ d.f. Degrees of freedom $= n-1$ because $\sum(X_i - \bar{\mathbf{X}}) = \mathbf{0}$ removes one.
  \item \textbf{$p = 1$ special case.} $W_1(m, \sigma^2) = \sigma^2 \chi^2_m$. The chi-squared is a special case of the Wishart.
  \item \textbf{Independence.} $\bar{\mathbf{X}}$ and $\mathbf{S}$ are independent (for Gaussian data). This is what makes Hotelling's $T^2$ test possible.
  \item \textbf{Expected value.} $\mathbb{E}[W] = m\boldsymbol{\Sigma}$ $\Rightarrow$ $\mathbb{E}[\mathbf{S}] = \boldsymbol{\Sigma}$ (unbiased), $\mathbb{E}[\hat{\boldsymbol{\Sigma}}] = \frac{n-1}{n}\boldsymbol{\Sigma}$ (biased).
  \item \textbf{Additivity.} $W_1 \sim W_p(m_1, \boldsymbol{\Sigma})$, $W_2 \sim W_p(m_2, \boldsymbol{\Sigma})$ independent $\Rightarrow$ $W_1 + W_2 \sim W_p(m_1 + m_2, \boldsymbol{\Sigma})$.
  \item \textbf{Positive definiteness.} $W \sim W_p(m, \boldsymbol{\Sigma})$ is almost surely positive definite $\Leftrightarrow m \geq p$. If $n \leq p$: $\mathbf{S}$ is singular, $\mathbf{S}^{-1}$ does not exist, standard tests fail.
  \item \textbf{Scalar projection.} $\mathbf{a}^\top W \mathbf{a} \sim (\mathbf{a}^\top \boldsymbol{\Sigma} \mathbf{a})\,\chi^2_m$ for any fixed $\mathbf{a} \neq \mathbf{0}$.
\end{itemize}

\subsection*{LLN and CLT}
\begin{itemize}[itemsep=5pt]
  \item \textbf{LLN (weak).} $\bar{\mathbf{X}}_n \xrightarrow{P} \boldsymbol{\mu}$ as $n \to \infty$, for any iid distribution with finite mean and covariance. No Gaussianity required.
  \item \textbf{LLN for covariance.} $\mathbf{S}_n \xrightarrow{P} \boldsymbol{\Sigma}$. Both estimators are \textbf{consistent}.
  \item \textbf{CLT.} $\sqrt{n}(\bar{\mathbf{X}}_n - \boldsymbol{\mu}) \xrightarrow{d} N_p(\mathbf{0}, \boldsymbol{\Sigma})$ for any distribution with finite $\boldsymbol{\Sigma}$. The $\sqrt{n}$ rescaling is mandatory --- without it the error collapses to zero.
  \item \textbf{Rate of convergence.} $\|\bar{\mathbf{X}}_n - \boldsymbol{\mu}\| \sim C/\sqrt{n}$. Halving error $\Rightarrow$ 4$\times$ data. Error$/10 \Rightarrow 100\times$ data.
  \item \textbf{Exam rule.} ``Gaussian data'' $\Rightarrow$ use exact Wishart/$\chi^2$ results. ``Large sample'' or ``general distribution'' $\Rightarrow$ use CLT. Exact results are always stronger.
  \item \textbf{The $n > p$ barrier.} If $n \leq p$, $\mathbf{S}$ is singular: you cannot compute $\mathbf{S}^{-1}$, and all classical tests break down. In practice $n \gg p$ is required.
  \item \textbf{Why is the Gaussian universal?} Not because data is Gaussian, but because the CLT forces \textit{averages of anything} toward Gaussian. It is the universal attractor of sums.
\end{itemize}

\subsection*{The Complete Scalar--Matrix Analogy (memorise this table)}
\renewcommand{\arraystretch}{1.4}
\begin{center}
\begin{tabular}{ll}
\hline
\textbf{Scalar ($p=1$)} & \textbf{Matrix ($p \geq 2$)} \\
\hline
$X_i \sim N(\mu, \sigma^2)$ & $X_i \sim N_p(\boldsymbol{\mu}, \boldsymbol{\Sigma})$ \\
$\bar{X} \sim N(\mu,\, \sigma^2/n)$ & $\bar{\mathbf{X}} \sim N_p(\boldsymbol{\mu},\, \boldsymbol{\Sigma}/n)$ \\
$(n-1)S^2/\sigma^2 \sim \chi^2_{n-1}$ & $(n-1)\mathbf{S} \sim W_p(n-1, \boldsymbol{\Sigma})$ \\
$\bar{X}$ and $S^2$ independent & $\bar{\mathbf{X}}$ and $\mathbf{S}$ independent \\
$z$-test / $t$-test & Hotelling's $T^2$ \\
$\chi^2_m = \sum_{j=1}^m Z_j^2$ & $W_p(m,\boldsymbol{\Sigma}) = \sum_{j=1}^m Z_j Z_j^\top$ \\
\hline
\end{tabular}
\end{center}


\end{document}